\documentclass[11pt]{article}
\usepackage[margin=1in]{geometry}
\usepackage{amsmath,amssymb,mathtools,bm}

\newcommand{\Exp}{\operatorname{Exp}}
\newcommand{\SO}{\mathrm{SO}}
\newcommand{\R}{\mathbb{R}}

\begin{document}

\section{Continuous-Time ESKF Error-State Derivation}

\subsection{Nominal state and IMU motion model}

Define the nominal state
\[
x=
\begin{bmatrix}
p & v & R & b_g & b_a & g
\end{bmatrix}^{T}.
\]

The continuous-time IMU motion equations are
\begin{align}
\dot p &= v,\\
\dot v &= R(\tilde a-b_a-\eta_a)+g,\\
\dot R &= R(\tilde\omega-b_g-\eta_g)^\wedge,\\
\dot b_g &= \eta_{bg},\\
\dot b_a &= \eta_{ba},\\
\dot g &= 0.
\end{align}

The IMU measurement equations may be written as
\begin{align}
\tilde a &= R^T(a-g)+b_a+\eta_a,\\
\tilde\omega &= \omega+b_g+\eta_g,
\end{align}
together with
\[
\dot R=R\omega^\wedge.
\]

For the nominal state, the kinematic equations have the same form but
ignore the measurement-noise terms:
\begin{align}
\dot p &= v,\\
\dot v &= R(\tilde a-b_a)+g,\\
\dot R &= R(\tilde\omega-b_g)^\wedge,\\
\dot b_g &= 0,\\
\dot b_a &= 0,\\
\dot g &= 0.
\end{align}

\subsection{True state}

Define the true state as
\[
x_t=
\begin{bmatrix}
p_t & v_t & R_t & b_{at} & b_{gt} & g_t
\end{bmatrix}^{T}.
\]

Its continuous-time dynamics are
\begin{align}
\dot p_t &= v_t,\\
\dot v_t &= R_t(\tilde a-b_{at}-\eta_a)+g_t,\\
\dot R_t &= R_t(\tilde\omega-b_{gt}-\eta_g)^\wedge,\\
\dot b_{gt} &= \eta_{bg},\\
\dot b_{at} &= \eta_{ba},\\
\dot g_t &= 0.
\end{align}

\subsection{Definition of the error state}

Use a right-multiplicative rotation error:
\begin{align}
p_t &= p+\delta p,\\
v_t &= v+\delta v,\\
R_t &= R\,\delta R
    =R\Exp(\delta\theta),\\
b_{gt} &= b_g+\delta b_g,\\
b_{at} &= b_a+\delta b_a,\\
g_t &= g+\delta g.
\end{align}

The quantities without the subscript $t$ are nominal-state variables.
The goal is to derive the continuous-time dynamics of
\[
\delta x=
\begin{bmatrix}
\delta p &
\delta v &
\delta\theta &
\delta b_g &
\delta b_a &
\delta g
\end{bmatrix}^{T}.
\]

\subsection{Position error}

Starting from
\[
p_t=p+\delta p,
\]
differentiate with respect to time:
\[
\dot p_t=\dot p+\delta\dot p.
\]

Since
\[
\dot p_t=v_t=v+\delta v
\]
and
\[
\dot p=v,
\]
we obtain
\[
v+\delta v=v+\delta\dot p.
\]

Therefore
\[
\boxed{
\delta\dot p=\delta v.
}
\]

\subsection{Velocity error}

Start from
\[
v_t=v+\delta v.
\]

Differentiating,
\[
\dot v_t=\dot v+\delta\dot v.
\]

Using the true-state equation,
\[
\dot v_t
=
R_t(\tilde a-b_{at}-\eta_a)+g_t.
\]

Substitute the error-state definitions:
\[
\dot v_t
=
R\Exp(\delta\theta)
(\tilde a-b_a-\delta b_a-\eta_a)
+g+\delta g.
\]

For small $\delta\theta$,
\[
\Exp(\delta\theta)
=
I+\delta\theta^\wedge
+O(\|\delta\theta\|^2).
\]

Keeping only first-order error terms,
\begin{align}
\dot v_t
&\approx
R(I+\delta\theta^\wedge)
(\tilde a-b_a-\delta b_a-\eta_a)
+g+\delta g\\
&\approx
R\tilde a-Rb_a-R\delta b_a-R\eta_a
+R\delta\theta^\wedge\tilde a
-R\delta\theta^\wedge b_a
+g+\delta g.
\end{align}

Terms such as
\[
\delta\theta^\wedge\delta b_a,
\qquad
\delta\theta^\wedge\eta_a
\]
are neglected as higher-order terms in this first-order error-state
linearization.

Using
\[
a^\wedge b=-b^\wedge a,
\]
we have
\[
\delta\theta^\wedge\tilde a
=
-\tilde a^\wedge\delta\theta,
\qquad
-\delta\theta^\wedge b_a
=
b_a^\wedge\delta\theta.
\]

Hence
\[
\dot v_t
\approx
R(\tilde a-b_a)+g
-R(\tilde a-b_a)^\wedge\delta\theta
-R\delta b_a
-R\eta_a
+\delta g.
\]

On the other hand, the nominal equation gives
\[
\dot v+\delta\dot v
=
R(\tilde a-b_a)+g+\delta\dot v.
\]

Comparing the two expressions,
\[
\boxed{
\delta\dot v
=
-R(\tilde a-b_a)^\wedge\delta\theta
-R\delta b_a
-R\eta_a
+\delta g.
}
\]

\subsection{Rotation error}

Start from
\[
R_t=R\Exp(\delta\theta).
\]

Differentiating,
\[
\dot R_t
=
\dot R\Exp(\delta\theta)
+
R\frac{d}{dt}\Exp(\delta\theta).
\]

\subsubsection*{Exact derivative of the exponential map}

For a general time-varying
\[
\delta\theta(t)\in\R^3,
\]
the exact derivative is
\[
\boxed{
\frac{d}{dt}\Exp(\delta\theta)
=
\Exp(\delta\theta)
\left(
J_r(\delta\theta)\delta\dot\theta
\right)^\wedge,
}
\]
where $J_r(\delta\theta)$ is the right Jacobian of $\SO(3)$.

This follows from
\[
\Exp(\delta\theta+\delta\dot\theta\,dt)
\approx
\Exp(\delta\theta)
\Exp\!\left(
J_r(\delta\theta)\delta\dot\theta\,dt
\right).
\]

Since
\[
J_r(\delta\theta)
=
I-\frac12\delta\theta^\wedge
+O(\|\delta\theta\|^2),
\]
the first-order ESKF approximation is
\[
J_r(\delta\theta)\delta\dot\theta
\approx
\delta\dot\theta.
\]

Therefore, within the first-order error-state derivation,
\[
\boxed{
\frac{d}{dt}\Exp(\delta\theta)
\approx
\Exp(\delta\theta)\delta\dot\theta^\wedge.
}
\]

Using this first-order approximation,
\[
\dot R_t
\approx
\dot R\Exp(\delta\theta)
+
R\Exp(\delta\theta)\delta\dot\theta^\wedge.
\]

The nominal rotation dynamics are
\[
\dot R
=
R(\tilde\omega-b_g)^\wedge.
\]

Hence
\[
\dot R_t
\approx
R(\tilde\omega-b_g)^\wedge\Exp(\delta\theta)
+
R\Exp(\delta\theta)\delta\dot\theta^\wedge.
\]

The true rotation dynamics are
\[
\dot R_t
=
R_t(\tilde\omega-b_{gt}-\eta_g)^\wedge
=
R\Exp(\delta\theta)
(\tilde\omega-b_{gt}-\eta_g)^\wedge.
\]

Equating both expressions and canceling $R$ gives
\begin{align}
\Exp(\delta\theta)\delta\dot\theta^\wedge
&=
\Exp(\delta\theta)
(\tilde\omega-b_{gt}-\eta_g)^\wedge\\
&\quad
-(\tilde\omega-b_g)^\wedge
\Exp(\delta\theta).
\end{align}

Use the $\SO(3)$ identity
\[
\phi^\wedge R
=
R(R^T\phi)^\wedge.
\]

With
\[
R=\Exp(\delta\theta),
\]
this gives
\[
(\tilde\omega-b_g)^\wedge\Exp(\delta\theta)
=
\Exp(\delta\theta)
\left(
\Exp(-\delta\theta)(\tilde\omega-b_g)
\right)^\wedge.
\]

Therefore
\begin{align}
\Exp(\delta\theta)\delta\dot\theta^\wedge
&=
\Exp(\delta\theta)
\Big[
(\tilde\omega-b_{gt}-\eta_g)^\wedge\\
&\qquad
-\left(
\Exp(-\delta\theta)(\tilde\omega-b_g)
\right)^\wedge
\Big].
\end{align}

For small $\delta\theta$,
\[
\Exp(-\delta\theta)
\approx
I-\delta\theta^\wedge.
\]

Thus
\begin{align}
&(\tilde\omega-b_{gt}-\eta_g)
-\Exp(-\delta\theta)(\tilde\omega-b_g)\\
&\approx
(\tilde\omega-b_{gt}-\eta_g)
-(I-\delta\theta^\wedge)(\tilde\omega-b_g)\\
&=
b_g-b_{gt}-\eta_g
+\delta\theta^\wedge\tilde\omega
-\delta\theta^\wedge b_g.
\end{align}

Since
\[
b_{gt}=b_g+\delta b_g,
\]
we have
\[
b_g-b_{gt}=-\delta b_g.
\]

Also,
\[
\delta\theta^\wedge(\tilde\omega-b_g)
=
-(\tilde\omega-b_g)^\wedge\delta\theta.
\]

Therefore
\[
\delta\dot\theta
=
-(\tilde\omega-b_g)^\wedge\delta\theta
-\delta b_g
-\eta_g.
\]

Hence
\[
\boxed{
\delta\dot\theta
=
-(\tilde\omega-b_g)^\wedge\delta\theta
-\delta b_g
-\eta_g.
}
\]

\subsubsection*{Identity used above}

For $R\in\SO(3)$,
\[
\operatorname{Ad}_R\phi
=
(R\phi^\wedge R^T)^\vee
=
R\phi.
\]

Thus
\[
R\phi^\wedge R^T=(R\phi)^\wedge.
\]

From this,
\[
\phi^\wedge R
=
R(R^T\phi)^\wedge.
\]

\subsection{Gyroscope-bias error}

Starting from
\[
b_{gt}=b_g+\delta b_g,
\]
differentiate:
\[
\dot b_{gt}
=
\dot b_g+\delta\dot b_g.
\]

The true bias dynamics are
\[
\dot b_{gt}=\eta_{bg},
\]
while the nominal model ignores noise:
\[
\dot b_g=0.
\]

Therefore
\[
\boxed{
\delta\dot b_g=\eta_{bg}.
}
\]

\subsection{Accelerometer-bias error}

Similarly,
\[
b_{at}=b_a+\delta b_a.
\]

Differentiating,
\[
\dot b_{at}
=
\dot b_a+\delta\dot b_a.
\]

Using
\[
\dot b_{at}=\eta_{ba},
\qquad
\dot b_a=0,
\]
we obtain
\[
\boxed{
\delta\dot b_a=\eta_{ba}.
}
\]

\subsection{Gravity error}

Starting from
\[
g_t=g+\delta g,
\]
differentiate:
\[
\dot g_t=\dot g+\delta\dot g.
\]

Since
\[
\dot g_t=0,
\qquad
\dot g=0,
\]
we obtain
\[
\boxed{
\delta\dot g=0.
}
\]

\section{Final continuous-time error-state equations}

Collecting the results,
\[
\boxed{
\begin{aligned}
\delta\dot p
&=
\delta v,\\[1mm]
\delta\dot v
&=
-R(\tilde a-b_a)^\wedge\delta\theta
-R\delta b_a
-R\eta_a
+\delta g,\\[1mm]
\delta\dot\theta
&=
-(\tilde\omega-b_g)^\wedge\delta\theta
-\delta b_g
-\eta_g,\\[1mm]
\delta\dot b_g
&=
\eta_{bg},\\[1mm]
\delta\dot b_a
&=
\eta_{ba},\\[1mm]
\delta\dot g
&=
0.
\end{aligned}
}
\]

\end{document}
